% =====================================================================
% Preprint (español) — diseño in silico de ASO para ABCA4 c.161-395G>A
%
% Compilar:  tectonic docs/preprint_es/preprint_es.tex
% Figuras:   scripts/run-in-env.sh python pipeline/make_figures.py
%            Se generan desde data/results/, así que texto y figura no pueden
%            divergir. Las etiquetas están en español (ver pipeline/make_figures.py).
%
% Versión en inglés: docs/preprint_en/preprint_en.tex — mismo contenido.
% =====================================================================
\documentclass[11pt,a4paper]{article}

\usepackage[spanish,es-tabla,es-noquoting]{babel}
\usepackage[utf8]{inputenc}
\usepackage[T1]{fontenc}
\usepackage{geometry}
\geometry{margin=2.5cm}

\usepackage{microtype}
\usepackage{graphicx}
\usepackage{booktabs}
\usepackage{array}
\usepackage{amsmath}
\usepackage{amssymb}
\usepackage{bm}
\usepackage{enumitem}
\usepackage{caption}
\usepackage{float}
\usepackage[section]{placeins}
\usepackage{needspace}
\usepackage{xcolor}
\usepackage{fancyhdr}
\usepackage{titlesec}
\usepackage[hidelinks]{hyperref}

\hypersetup{colorlinks=true, linkcolor=blue!50!black, citecolor=blue!50!black,
            urlcolor=blue!50!black,
            pdftitle={Un pipeline construido por un agente para diseno in silico de ASO},
            pdfauthor={Sergio Mauricio Nunez; Amyra Sanchez}}

% --- composición: que ningún bloque se parta mal ----------------------
\widowpenalty=10000
\clubpenalty=10000
\displaywidowpenalty=10000
\brokenpenalty=10000
\raggedbottom
\sloppy
\emergencystretch=3em
\setlength{\parskip}{0.45em}
\setlength{\parindent}{0pt}

\newcommand{\brk}{\discretionary{}{}{}}

% Los signos ¿ y ¡ literales salen mal con este motor (tectonic/XeTeX): U+00BF
% cae en un slot equivocado de la fuente T1 y se imprime «£». Con las macros
% salen bien. Verificado el 2026-08-02; si algún día se compila con pdflatex,
% esto sigue siendo correcto, solo que innecesario.
\newcommand{\qd}{\textquestiondown}
\newcommand{\ed}{\textexclamdown}

\titleformat{\section}{\large\bfseries}{\thesection}{0.8em}{}
\titleformat{\subsection}{\normalsize\bfseries}{\thesubsection}{0.8em}{}
\titlespacing*{\section}{0pt}{1.4em}{0.5em}
\titlespacing*{\subsection}{0pt}{1.0em}{0.35em}

\captionsetup{font=small, labelfont=bf, width=0.94\linewidth}

\pagestyle{fancy}
\fancyhf{}
\fancyhead[L]{\small Diseño in silico de ASO para \emph{ABCA4} c.161-395G>A}
\fancyhead[R]{\small Preprint}
\fancyfoot[C]{\thepage}
\renewcommand{\headrulewidth}{0.4pt}

\begin{document}

% =====================================================================
% Portada — armada a mano en vez de \maketitle, que no lograba mantener
% el bloque de autores dentro del ancho de texto.
% =====================================================================
\thispagestyle{empty}
\begin{center}
  \vspace*{0.6cm}

  {\LARGE\bfseries Un pipeline construido por un agente para\\[0.15em]
   el diseño in silico de oligonucleótidos antisentido,\\[0.25em]
   y lo que seis revisores adversariales encontraron en él\par}

  \vspace{0.8em}
  {\large Informe metodológico sobre \emph{ABCA4} c.161-395G>A\par}

  \vspace{1.6em}
  {\large Sergio Mauricio Nu\~nez\qquad Amyra Sanchez\par}
  \vspace{0.4em}
  {\small \emph{YAIS Lab}\par}
  \vspace{0.2em}
  {\small \href{mailto:mauricio.nunez@yaislab.org}{mauricio.nunez@yaislab.org}\par}

  \vspace{0.9em}
  {\small 2 de agosto de 2026 \quad\textbullet\quad Preprint, sin revisión por pares\par}
\end{center}

\vspace{0.8em}

\begin{center}
\begin{minipage}{0.88\linewidth}\small\itshape
La implementación, la ejecución del pipeline y la redacción estuvieron a cargo de un
agente de IA (Claude, Anthropic) bajo supervisión de los autores. La
sección~\ref{sec:procedencia} declara explícitamente qué decisiones fueron de los
autores y cuáles no.
\end{minipage}
\end{center}

\vspace{1.2em}

\begin{center}
\begin{minipage}{0.92\linewidth}
{\bfseries\large Resumen}\\[0.4em]
\small
La variante intrónica profunda \textbf{\emph{ABCA4} c.161-395G>A} activa un pseudoexón aberrante
(PE1b, 91~pb) y fue el alelo intrónico profundo más frecuente en una cohorte china de enfermedad de
Stargardt publicada~\cite{wang2025}. No hay ningún oligonucleótido antisentido (ASO) publicado ni
patentado contra ella. Presentamos un pipeline computacional reproducible de siete módulos para el
diseño in silico de candidatos ASO de conmutación de splicing con química \textbf{PMO}~\cite{summerton1997},
construido de punta a punta por un agente de IA en seis días y sometido después a un panel de seis
revisores adversariales independientes.

El pipeline reproduce tres resultados con anclaje externo. Primero, la variante aporta la adenina de
consenso en $-2$ del motivo donador sobre un dinucleótido GT intacto
(\texttt{GGG|GTAGGT}~$\rightarrow$~\texttt{GAG|GTAGGT}), una observación sin modelo y sin umbral.
Segundo, dos predictores neuronales entrenados de forma independiente~\cite{jaganathan2019,zeng2022}
ubican su máximo exactamente en $+1$ y $-89$ respecto de la variante, lo que implica un pseudoexón
de 91~pb --- y eso reproduce, \emph{posición por posición}, los bordes que Corradi et
al.~\cite{corradi2023} ya habían medido por ensayo de midigén en 2023 y cuyo tamaño Wang et
al.~\cite{wang2025} confirmaría después por minigén. Tercero, sobre un oligo-walk publicado
independiente con eficacia medida~\cite{kaltak2023}, el mismo puntaje ubica los cinco AONs eficaces
conocidos en las posiciones 1, 3, 4, 5 y 6 de 32 ($\mathrm{AUC}=0{,}974$).

Reportamos la revisión adversarial como resultado primario y no como un párrafo de limitaciones. El
panel devolvió \emph{rechazo con invitación a resubmisión mayor} e identificó siete problemas
críticos. Corregirlos cambió las conclusiones: un error de orientación de hebra en el cálculo de la
temperatura de fusión se había propagado cuatro módulos aguas arriba y había generado un hallazgo,
un registro de decisión y una conclusión publicada; un control negativo con el tejido
deliberadamente equivocado reprodujo exactamente la ``concordancia entre tres predictores'', lo que
muestra que esa concordancia no informa; y una afirmación repetida en seis documentos durante cuatro
días nunca se había verificado contra la literatura. Armar la lista de referencias de este
manuscrito --- el último paso antes de someterlo --- produjo dos más del mismo tipo: el paper fuente
había estado atribuido al primer autor equivocado durante cinco días, y las coordenadas del
pseudoexón que el proyecto trataba como propias estaban publicadas desde 2023.

\textbf{Ningún candidato fue sintetizado ni ensayado.} Sostenemos que, para investigación construida
por agentes, la revisión adversarial no es control de calidad opcional sino parte del método, y
cuantificamos por qué.
\end{minipage}
\end{center}

\clearpage

% =====================================================================
\tableofcontents
\clearpage

% =====================================================================
\section{Motivación}

La enfermedad de Stargardt tipo~1 es la distrofia macular hereditaria más común, causada por
variantes bialélicas en \emph{ABCA4}. Una fracción relevante de los alelos patogénicos no son
variantes codificantes sino \textbf{variantes intrónicas profundas} que activan sitios de splicing
crípticos, llevando al espliceosoma a tratar un tramo intrónico como si fuera un exón --- un
\emph{pseudoexón} --- e introducir un codón de terminación prematuro.

La variante \textbf{c.161-395G>A} fue reportada por tres grupos independientes. Corradi et
al.~\cite{corradi2023} la caracterizaron primero, en 2023, por ensayo de midigén en una cohorte
europea, con 36\,\% de ARN silvestre residual y los bordes exactos del pseudoexón. Wang et
al.~\cite{wang2025} mostraron por minigén que activa tres pseudoexones (PE1b 91~pb, PE1c 255~pb,
PE1d 265~pb), con el 76,4\,\% del ARNm resultante incorrectamente spliceado, y que es el alelo
intrónico profundo más frecuente de su cohorte china. Huang et al.~\cite{huang2025} la recuperaron
otra vez en una cohorte taiwanesa. Ninguno de los tres reporta un oligonucleótido antisentido contra
ella.

Los ASO de conmutación de splicing se unen al pre-ARNm por complementariedad de bases y
\textbf{bloquean estéricamente} el acceso a un sitio de splicing sin degradar el
transcrito~\cite{dhuri2020}. Ningún ASO publicado o patentado apunta a esta variante --- lo que
motiva el trabajo, pero también significa que \textbf{no hay control positivo} con el cual validar
un pipeline de diseño de punta a punta. La sección~\ref{sec:calibracion} describe el sustituto
parcial que usamos.

% =====================================================================
\section{Procedencia y división del juicio}
\label{sec:procedencia}

Lo declaramos antes de los métodos porque condiciona cómo debe leerse el resto.

La implementación, la ejecución del pipeline y la redacción las produjo un agente de IA. Las
siguientes fueron decisiones de los autores, no del agente:

\begin{itemize}[nosep, leftmargin=1.4em]
  \item \textbf{La química.} PMO se eligió como decisión de alcance, en contra del precedente de la
        literatura retinal (2$'$-MOE/PS), y quedó documentada con sus costos.
  \item \textbf{Encargar la revisión adversarial}, con el razonamiento explícito de que un agente no
        puede auditar el trabajo que produjo.
  \item \textbf{Cuestionar las afirmaciones del agente.} Tres de los errores de mayor impacto
        salieron a la luz porque los autores preguntaron, no porque el agente los detectara --- igual
        que el hecho de que los parámetros de campo de fuerza para PMO sí
        existen~\cite{maksudov2023}, después de que el agente afirmara que no.
\end{itemize}

Consideramos esta declaración un requisito metodológico y no un agradecimiento.

% =====================================================================
\section{El pipeline}
\label{sec:pipeline}

Siete módulos, cada uno con su runner reproducible y sus limitaciones declaradas
(Tabla~\ref{tab:modulos}, Figura~\ref{fig:embudo}).

\begin{table}[htbp]
\centering
\small
\setlength{\tabcolsep}{4.5pt}
\begin{tabular}{@{}clll@{}}
\toprule
\textbf{\#} & \textbf{Módulo} & \textbf{Pregunta} & \textbf{Resultado} \\
\midrule
1  & Secuencia          & \qd Dónde está la variante?            & Coordenada confirmada por dos vías~\cite{freeman2018} \\
2  & Oligo-walk         & \qd Qué ventanas?                      & 381 candidatos \\
3  & Filtros heurísticos& \qd Oligos viables?                    & 381 $\rightarrow$ 276 \\
4  & Termodinámica      & \qd Se unen bien, son accesibles?      & 276 $\rightarrow$ 78 \\
5  & Off-target         & \qd Se parecen a otros genes?          & 78 anotados por severidad \\
6  & Splicing neuronal  & \qd El sitio falso es real?            & Donador $+1$, aceptor $-89$ \\
6b & Enmascarado del ASO& \qd Cada parche lo apaga?              & 12 anulan el pseudoexón \\
7  & Ranking            & \qd Cuál sintetizar primero?           & Frente de Pareto: 3 \\
\bottomrule
\end{tabular}
\caption{Los siete módulos y qué responde cada uno. Cada fila tiene un comando de regeneración en el
README del repositorio.}
\label{tab:modulos}
\end{table}

\begin{figure}[htbp]
\centering
\includegraphics[width=0.9\linewidth]{figs/fig1_funnel.png}
\caption{Embudo de candidatos con los números vigentes. Las dos etapas verdes son las que requieren
inferencia neuronal; todo lo anterior es aritmética de secuencia. El ancho de cada banda escala como
$n^{0{,}38}$ para que las últimas etapas sigan siendo visibles.}
\label{fig:embudo}
\end{figure}

\subsection{Cómo se le representa un ASO a un predictor de splicing}
\label{sec:enmascarado}

Los predictores de splicing toman una sola hebra de secuencia y no tienen representación de un
ligando unido. Un PMO no corta ni modifica el ARN --- lo tapa. Por eso reemplazamos la ventana del
ASO por \texttt{N}, que la codificación one-hot de ambas familias de predictores mapea al vector
nulo: ``acá no hay información de secuencia legible''. Es un proxy directo del bloqueo estérico.

Sea $x$ la secuencia del pre-ARNm mutante, $c$ una ventana candidata y $\mathcal{M}(x,c)$ la misma
secuencia con cada base dentro de $c$ reemplazada por \texttt{N}. Para un predictor $p$ y un sitio de
splicing $s$, escribimos $S_p(\cdot,s)$ para el score que asigna en esa posición. La cantidad que el
pipeline usa efectivamente es la \emph{fracción de señal retenida}:

\begin{equation}
r_p(c,s) \;=\; \frac{S_p\bigl(\mathcal{M}(x,c),\,s\bigr)}{S_p(x,s)}
\label{eq:retencion}
\end{equation}

$r_p = 1$ significa que la máscara no cambió nada; $r_p = 0$, que el sitio desapareció. Expresar el
efecto como cociente y no como diferencia es lo que hace al criterio comparable entre predictores
cuyas escalas de score difieren.

Se considera que un candidato \textbf{anula el pseudoexón} cuando apaga \emph{al menos uno} de sus
dos bordes --- el donador críptico $d$ en $+1$ o el aceptor críptico $a$ en $-89$ --- en
\emph{todos} los predictores:

\begin{equation}
V(c) = \text{anula} \iff \forall p \in P:\;
\min_{s \in \{d,\,a\}} r_p(c,s) \;\le\; \tau,
\qquad \tau = 0{,}25
\label{eq:veredicto}
\end{equation}

La disyunción entre bordes es una afirmación biológica y no una comodidad: un pseudoexón necesita que
se reconozcan sus dos bordes, así que quitar cualquiera de los dos alcanza para desarmarlo. El umbral
$\tau$ es una decisión de diseño declarada y no está calibrada contra experimento.

Se descartaron dos alternativas: mutar la ventana introduce motivos espurios, y recortarla desplaza
todas las coordenadas aguas abajo. Cabe notar que la propia rutina de inferencia de los predictores
ya rellena con \texttt{"N"}$\times$5000 a cada lado, así que los bloques de \texttt{N} son la
convención operativa de estos modelos y no una entrada inventada.

\subsection{Termodinámica, y el punto exacto donde la química no encaja}
\label{sec:termo}

El Módulo~4 estima la temperatura de fusión del dúplex ASO:ARN con el modelo de vecino más
cercano~\cite{santalucia1998}, usando la tabla de híbridos ARN/ADN de Sugimoto et
al.~\cite{sugimoto1995} tal como está implementada en Biopython~\cite{cock2009}:

\begin{equation}
T_m \;=\; \frac{\Delta H^\circ}{\;\Delta S^\circ + R\,\ln\!\bigl(C_T/x\bigr)\;} \;-\; 273{,}15
\label{eq:tm}
\end{equation}

con la entalpía y la entropía acumuladas sobre pasos dinucleotídicos consecutivos de la secuencia
$\bm{s}=s_1s_2\ldots s_n$,

\begin{equation}
\Delta H^\circ = \Delta H^\circ_{\mathrm{init}} + \sum_{i=1}^{n-1} \Delta H^\circ(s_i s_{i+1}),
\qquad
\Delta S^\circ = \Delta S^\circ_{\mathrm{init}} + \sum_{i=1}^{n-1} \Delta S^\circ(s_i s_{i+1}).
\label{eq:nn}
\end{equation}

\needspace{5\baselineskip}
\textbf{La tabla es asimétrica, y eso es lo que se rompió.} Para un dúplex híbrido el juego de
parámetros está indexado por la hebra de \textbf{ARN}, de modo que
$\Delta H^\circ(\mathtt{AG/TT}) \neq \Delta H^\circ(\mathtt{TT/AG})$. En los valores tabulados estos
difieren en más de un factor de dos ($-3{,}96$ frente a $-10{,}38$~kcal\,mol$^{-1}$; $-11{,}6$
frente a $-31{,}8$~cal\,mol$^{-1}$K$^{-1}$). Pasar el ASO donde se esperaba la diana de ARN no
produce entonces una perturbación pequeña --- produce otro número. La
sección~\ref{sec:bugtm} reporta lo que eso costó.

\textbf{Dos sesgos que no se cancelan.} La ecuación~\ref{eq:tm} supone un esqueleto cargado, y la
corrección salina que se le aplica encima supone $\sim$50~mM de Na$^+$. El esqueleto de un PMO es
neutro~\cite{summerton1997}, y las mediciones publicadas de fusión de PMO se hacen cerca de 1,0~M de
NaCl. De forma independiente, se reporta que los enlaces fosforodiamidato \emph{estabilizan} el
dúplex~\cite{maksudov2023}. Los dos efectos empujan la estimación en el mismo sentido: el proxy
\textbf{subestima} la temperatura de fusión real. Por eso hoy $T_m$ se reporta como anotación y no
se usa como filtro (sección~\ref{sec:recuperado}).

\subsection{Off-target}

El Módulo~5 busca cada candidato contra el transcriptoma humano con BLAST+~\cite{camacho2009} y
registra $L(c)$, el tramo contiguo perfecto más largo contra cualquier transcrito que no sea
\emph{ABCA4}. La severidad se anota, no se usa como gate: para una química de bloqueo estérico, a
diferencia de un gapmer que recluta RNasa~H, la homología solo importa según dónde
caiga~\cite{dhuri2020}. Dos resultados publicados acotan cuánto puede confiarse en esta medición:
ningún rasgo de secuencia por sí solo predice bien el mis-splicing off-target real, y el propio
BLAST penaliza gaps e ignora el apareamiento oscilante G:U~\cite{scharner2020}; y el número de
sitios complementarios crece abruptamente con la tolerancia a mismatches que se
permita~\cite{yoshida2018}.

\subsection{Ranking: por qué un frente de Pareto y no un puntaje ponderado}
\label{sec:pareto}

Las tres dimensiones del ranking tienen unidades incomparables. Todas se escriben de modo que ``más
alto es mejor'':

\begin{align}
f_1(c) &= 1 - \frac{1}{|P|}\sum_{p \in P} \min_{s \in \{d,a\}} r_p(c,s)
        && \text{fuerza de bloqueo} \label{eq:f1}\\[0.2em]
f_2(c) &= -\,L(c)
        && \text{seguridad off-target} \label{eq:f2}\\[0.2em]
f_3(c) &= \tfrac{1}{2}\bigl(P_{\mathrm{acc}}(c) + P_{\mathrm{hd}}(c)\bigr)
        && \text{calidad termodinámica} \label{eq:f3}
\end{align}

donde $P_{\mathrm{acc}}$ y $P_{\mathrm{hd}}$ son los percentiles dentro de la cohorte de la
accesibilidad de la ventana diana y de la no-autodimerización, calculados con
ViennaRNA~\cite{lorenz2011}.

Promediar $f_1$, $f_2$ y $f_3$ en un solo puntaje exige postular una tasa de cambio --- cuántos pares
de bases de homología off-target valen 0,01 de fuerza de bloqueo --- que \textbf{ningún dato de este
proyecto respalda}, precisamente porque no hay control positivo contra el cual calibrarla. Un frente
de Pareto no necesita esa tasa. Escribiendo $\bm{f}(c) = (f_1,f_2,f_3)$, el candidato $a$
\emph{domina} a $b$ cuando

\begin{equation}
a \succ b
\iff
\bigl(\forall i:\; f_i(a) \ge f_i(b)\bigr)
\;\wedge\;
\bigl(\exists j:\; f_j(a) > f_j(b)\bigr),
\label{eq:dominancia}
\end{equation}

y el frente es el conjunto de candidatos a los que nadie domina, restringido a los que pasan la
puerta de elegibilidad $V(c)=\text{anula}$ de la ecuación~\ref{eq:veredicto}:

\begin{equation}
\mathcal{F} \;=\; \bigl\{\, a \in E \;:\; \nexists\, b \in E,\; b \succ a \,\bigr\},
\qquad
E = \{\,c : V(c) = \text{anula}\,\}.
\label{eq:frente}
\end{equation}

Devuelve un conjunto y no un ganador, que es la salida honesta.

\needspace{6\baselineskip}
\textbf{La convención que hay que declarar.} Colapsar accesibilidad y no-autodimerización en un solo
promedio en la ecuación~\ref{eq:f3} \emph{es} una ponderación implícita 50/50 --- exactamente lo que
esta sección dice evitar. Lo hacemos igual, por una razón medida y no estética: manteniéndolas como
dimensiones separadas, el frente pasa de 3 a \textbf{11 de los 12} candidatos elegibles, y un frente
que casi no descarta nada casi no informa nada. El código reporta ese número para que la decisión
quede a la vista en lugar de enterrada.

% =====================================================================
\section{Calibración contra un oligo-walk independiente}
\label{sec:calibracion}

Como no existe ningún ASO para esta variante, el pipeline se aplicó sin cambios a otra variante de
\emph{ABCA4} que \emph{sí} tiene eficacia medida: c.5461-10T>C en el intrón~38, para la cual Kaltak
et al.~\cite{kaltak2023} publicaron un walk de 32 AONs y reportaron cinco como eficaces, uno de los
cuales se convirtió en el candidato clínico QR-1011. El módulo de calibración es deliberadamente
paralelo al pipeline principal y no lo toca.

Puntuar los 32 AONs con el mismo $\Delta$ de enmascarado de la sección~\ref{sec:enmascarado} y
preguntar dónde caen los eficaces conocidos da un estadístico de Mann--Whitney. Con $K$ el conjunto
de eficaces conocidos y $O$ el resto,

\begin{equation}
U = \sum_{k \in K} \sum_{o \in O}
\Bigl[\mathbf{1}\bigl(\delta_k > \delta_o\bigr) + \tfrac{1}{2}\,\mathbf{1}\bigl(\delta_k = \delta_o\bigr)\Bigr],
\qquad
\mathrm{AUC} = \frac{U}{|K|\,|O|}.
\label{eq:auc}
\end{equation}

Los cinco AONs eficaces conocidos ocupan las posiciones \textbf{1, 3, 4, 5 y 6 de 32}, con QR-1011
tercero: $\mathrm{AUC} = 0{,}974$. Repetir toda la calibración con el predictor entrenado en
retina~\cite{riepe2025} da las mismas cinco posiciones y $\mathrm{AUC} = 0{,}970$, así que la
separación no es un artefacto de un modelo. Un análisis de sensibilidad que quita los cuatro AONs que
cubren el aceptor del exón~39 --- donde anular el sitio es trivialmente malo --- deja 28 AONs y
$\mathrm{AUC} > 0{,}95$.

\textbf{Qué habilita y qué no.} Muestra que el puntaje no es arbitrario: discrimina AONs con eficacia
medida sobre una diana para la que nunca fue ajustado. No se transfiere a c.161-395G>A, que es otro
intrón, otro mecanismo (inclusión de pseudoexón en vez de salto de exón) y otra química.

% =====================================================================
\section{Resultados con anclaje externo}

Estos dos no dependen de ninguna métrica definida por nosotros.

\subsection{La variante aporta la adenina de consenso del donador}

Contando coincidencias contra el consenso donador \texttt{MAG|GTRAGT}, la secuencia silvestre 1~nt
aguas abajo de la variante lee \texttt{GGG|GTAGGT} y la mutante lee \texttt{GAG|GTAGGT}. La variante
misma aporta la adenina en la posición $-2$ --- una de las posiciones más conservadas del consenso
donador --- sobre un dinucleótido GT intacto, con purinas en $+3$ y $+5$. Esto no requiere modelo, ni
umbral, ni predictor.

\subsection{El pseudoexón predicho reproduce coordenadas publicadas, posición por posición}
\label{sec:coordenadas}

La inferencia local de SpliceAI~\cite{jaganathan2019} sobre una ventana de 10~kb predice un donador
críptico reforzado en $+1$ ($0{,}194 \rightarrow 0{,}560$) y un aceptor críptico en $-89$
($0{,}061 \rightarrow 0{,}271$), lo que implica un pseudoexón de \textbf{91~pb}.

Corradi et al.~\cite{corradi2023} midieron esta variante con un ensayo de midigén en 2023 y
reportaron el transcrito resultante como \texttt{r.160\_161ins161-484\_161-394}. Convertido a
offsets relativos a la variante, que está en \texttt{c.161-395}, esa variante de ARN publicada dice
que el aceptor está en $395-484 = -89$, el donador en $395-394 = +1$, y el pseudoexón mide
$484-394+1 = 91$~nt. \textbf{La predicción coincide con el experimento publicado en los dos bordes,
no solo en la longitud.} Wang et al.~\cite{wang2025} confirmaron después el tamaño de 91~pb por
minigén, como PE1b.

Una tercera cohorte aporta una replicación de otro tipo. Huang et al.~\cite{huang2025} reportan
delta scores de SpliceAI de $\mathrm{DS\_AG} = 0{,}21$ y $\mathrm{DS\_DG} = 0{,}37$ para esta
variante; nuestra inferencia local da $0{,}271-0{,}061 = 0{,}210$ y $0{,}560-0{,}194 = 0{,}366$.
Otro grupo, otro pipeline, otra cohorte, los mismos deltas a dos decimales.

Pangolin~\cite{zeng2022}, una arquitectura distinta entrenada por otro grupo, ubica su máximo en
\emph{exactamente} las mismas dos posiciones, con las posiciones vecinas en $\sim$1e$-$5. Un intento
deliberado de romper este resultado durante la revisión adversarial falló.

\needspace{5\baselineskip}
\textbf{Esto es una validación, no un descubrimiento.} Lo declaramos sin rodeos: las coordenadas ya
estaban en la literatura. Corradi et al.\ las publicaron dos años antes de que este proyecto
empezara, y el proyecto no lo supo hasta que se armó la lista de referencias de este preprint
(sección~\ref{sec:novedad}). Lo que la coincidencia establece es que el pipeline recupera desde
secuencia sola un pseudoexón medido funcionalmente --- que es para lo que hace falta confiar en un
pipeline de diseño, y que acá vale más de lo que habría valido una afirmación de novedad.

\textbf{La variante refuerza un sitio que ya existía.} Nuestro propio score del donador silvestre es
$0{,}194$, que no es cero, y Huang et al.~\cite{huang2025} consultaron SpliceVault y encontraron que
los dos sitios crípticos se usan en transcriptomas normales, por debajo del 3\,\% de los
transcritos. La formulación correcta es entonces que el cambio de base \emph{refuerza} un sitio
críptico débil preexistente --- consistente con la observación de la adenina de consenso de más
arriba --- y no que lo crea \emph{de novo}. Borradores anteriores de este trabajo decían ``crea'';
era demasiado fuerte.

\textbf{Lo que no se sigue.} SpliceAI no explica PE1c ni PE1d: sus aceptores implicados puntúan
0,016 y 0,042, muy por debajo de cualquier umbral utilizable. El modelo da cuenta de uno de los tres
pseudoexones reportados, y lo declaramos en vez de omitirlo.

\subsection{Qué ventanas apagan el pseudoexón}

Aplicar las ecuaciones~\ref{eq:retencion} y~\ref{eq:veredicto} a los 78 candidatos que llegan al
Módulo~6b da la Figura~\ref{fig:enmascarado}. Aparecen dos mecanismos distintos, y no son
intercambiables.

\begin{figure}[htbp]
\centering
\includegraphics[width=\linewidth]{figs/fig2_masking.png}
\caption{Efecto de enmascarar cada ventana candidata sobre la señal del donador críptico, como
fracción del basal sin máscara (ecuación~\ref{eq:retencion}). Las ventanas que cubren el donador en
$+1$ lo llevan a cero. Nueve candidatos centrados cerca de $-110$ anulan el \emph{aceptor} sin cubrir
el donador, y acá se ven solo como una caída parcial de la señal del donador --- lo que evalúa la
ecuación~\ref{eq:veredicto} es su retención del aceptor. La línea de trazos es el umbral declarado
$\tau = 0{,}25$.}
\label{fig:enmascarado}
\end{figure}

% =====================================================================
\section{La revisión adversarial como resultado}
\label{sec:revision}

Seis revisores independientes --- biología del splicing, estadística, reproducibilidad, terapéutica
de ASO, integridad del código y un equipo rojo hostil --- recibieron un dossier con cada afirmación,
su nivel de evidencia y su comando de regeneración, más reglas de juego explícitas que exigían que
cada objeción citara un archivo, una línea o un número.

El veredicto fue \emph{rechazo con invitación a resubmisión mayor}. Se plantearon siete problemas
críticos. Cuatro se verificaron de forma independiente antes de aceptarlos. La
Figura~\ref{fig:impacto} resume qué cambió; las subsecciones dan el mecanismo de cada uno.

\begin{figure}[htbp]
\centering
\includegraphics[width=\linewidth]{figs/fig3_review_impact.png}
\caption{Qué cambió la revisión. \textbf{Izquierda:} candidatos que sobreviven al Módulo~4 y, de
esos, cuántos anulan el pseudoexón. \textbf{Derecha:} candidatos que anulan el aceptor críptico sin
cubrirlo --- la observación que había motivado el criterio de veredicto del proyecto. Existía antes
de la revisión por la razón equivocada (un bug), desapareció por completo al corregirlo, y reapareció
sobre base firme cuando $T_m$ dejó de usarse como filtro.}
\label{fig:impacto}
\end{figure}

\subsection{Un bug del Módulo 4 que viajó cuatro módulos aguas arriba}
\label{sec:bugtm}

\texttt{Bio.SeqUtils.\brk MeltingTemp.\brk Tm\_NN} con la tabla \texttt{R\_DNA\_NN1} documenta que el
argumento debe ser la hebra de \textbf{ARN}; el código pasaba el ASO. Como se muestra en la
sección~\ref{sec:termo}, la tabla es asimétrica, así que el orden cambia el número.

Corregirlo movió el embudo del Módulo~4 de \textbf{44 a 16} candidatos, con solo 6 en común. Más
grave: los siete candidatos que anulaban el aceptor \emph{sin cubrirlo} --- la observación que había
motivado el criterio de veredicto del proyecto, su propio registro de decisión y una conclusión
publicada sobre concordancia entre predictores --- \textbf{no sobrevivieron}. Un defecto en el módulo
más barato había generado un hallazgo cuatro módulos más abajo: no un número equivocado, sino una
afirmación equivocada sobre biología.

\subsection{Un control de tejido equivocado que disolvió un argumento}

El proyecto había reportado que tres predictores entrenados de forma independiente, uno de ellos
específico de retina~\cite{riepe2025}, seleccionaban exactamente el mismo conjunto de candidatos, y
presentaba eso como validación cruzada. El panel propuso el control que faltaba: correr un modelo
entrenado \emph{deliberadamente} sobre el tejido equivocado --- el juego de pesos GTEx publicado
junto a Retina-SpliceAI, que excluye la retina neural (Figura~\ref{fig:tejido}).

\textbf{El tejido equivocado selecciona los candidatos del modelo de retina uno por uno}: ambos
devuelven los mismos 11 de 78, candidato por candidato, y no meramente la misma cantidad. La
concordancia refleja entonces arquitectura compartida y anotaciones de entrenamiento compartidas, no
biología robusta al tejido. El control costó seis minutos de CPU y estaba disponible desde que se
instalaron los pesos; nadie en el proyecto lo había corrido, porque nadie estaba buscando una forma
de refutar su propio argumento.

\begin{figure}[htbp]
\centering
\includegraphics[width=\linewidth]{figs/fig4_tissue_control.png}
\caption{El control negativo. \textbf{Izquierda:} los tres juegos de pesos --- agnóstico de tejido,
entrenado en retina, y un modelo GTEx que excluye deliberadamente la retina neural --- coinciden en
que la variante refuerza el donador críptico, aunque discrepan en su fuerza absoluta. \textbf{Derecha:}
el modelo de retina y el del tejido equivocado seleccionan conjuntos \emph{idénticos} de candidatos
(11 de 78); el agnóstico de tejido difiere en uno. Como el tejido equivocado reproduce el resultado,
la concordancia entre ellos no puede ser evidencia de biología robusta al tejido.}
\label{fig:tejido}
\end{figure}

\subsection{Una afirmación repetida cuatro días sin verificar}

La afirmación de que no existen parámetros termodinámicos de vecino más cercano para PMO se propagó
por seis documentos a partir de una sola inferencia no verificada. Comprobarla encontró: la
afirmación es esencialmente correcta para PMO, pero el modelo NN \emph{sí} está documentado como
aplicable a oligómeros sin carga; \emph{sí} existen parámetros para 2$'$-MOE/PS~\cite{park2025}, la
química de la literatura retinal real; y los parámetros de campo de fuerza para PMO \emph{sí}
existen~\cite{maksudov2023}. Además, la evidencia experimental muestra que los enlaces
fosforodiamidato \emph{estabilizan} el dúplex, así que el proxy \textbf{subestima} la temperatura de
fusión real --- un sesgo en la misma dirección que la corrección salina inaplicable, de modo que los
dos se acumulan en vez de cancelarse.

\subsection{Una corrección que recuperó lo que el filtro había escondido}
\label{sec:recuperado}

Como el umbral de temperatura de fusión no podía interpretarse para esta química, se lo convirtió de
filtro en anotación --- siguiendo un precedente que el proyecto ya había sentado para la severidad
off-target. El embudo pasó de 16 a \textbf{78} candidatos.

Los candidatos de solo-aceptor \textbf{reaparecieron}: nueve de ellos, con temperaturas de fusión
proxy de 36,5--40,1\,$^\circ$C, todas por debajo del umbral descartado. El mecanismo era real; el
cálculo con bug los había admitido por la razón equivocada y el gate corregido los había escondido
por completo. El frente de Pareto final vuelve a los mismos tres candidatos del análisis original,
ahora por una vía defendible.

\subsection{Una cita mal atribuida, encontrada al armar esta bibliografía}
\label{sec:atribucion}

Al compilar la lista de referencias de este preprint, la fuente de la medición por minigén se
verificó contra Crossref y PubMed Central por primera vez. El DOI y el identificador PMC que el
proyecto tenía registrados eran correctos; el nombre del autor no. El trabajo es
\textbf{Wang et al.}~\cite{wang2025} --- no hay ningún autor apellidado Peng en él. La atribución
estaba mal en 19 lugares del repositorio y en 15 documentos del vault, durante cinco días, y ya había
sido leída por los seis revisores sin que ninguno la señalara.

Lo reportamos porque es el mismo modo de falla que los otros tres: una afirmación que nunca se
verificó contra su fuente, repetida hasta que la repetición la hizo parecer verificada. Es también la
más barata de las cuatro de haber detectado --- una llamada a una API --- y, en ese momento, la
última que se encontró.

\subsection{Y las coordenadas que creíamos nuestras estaban publicadas desde 2023}
\label{sec:novedad}

El mismo trabajo de bibliografía produjo un quinto caso, y el de mayor consecuencia. El proyecto
tenía registrado que el paper fuente no reporta las posiciones de los sitios crípticos. Eso es
cierto de Wang et al.~\cite{wang2025} y se verificó contra su texto. Después el proyecto lo
generalizó a que \emph{nadie} las había reportado --- lo cual nunca se comprobó, y es falso.

Una búsqueda de texto completo de la literatura por la cadena de la variante devuelve exactamente
tres papers. El más antiguo, Corradi et al.~\cite{corradi2023}, midió la variante por ensayo de
midigén en septiembre de 2023 y publicó el pseudoexón con sus dos bordes, en la notación citada en
la sección~\ref{sec:coordenadas}. El tercero, Huang et al.~\cite{huang2025}, apareció en octubre de
2025 con una cohorte taiwanesa y reporta los mismos deltas de SpliceAI que nosotros calculamos.

De ahí se siguen dos cosas, y apuntan en direcciones opuestas. El anclaje externo es
\textbf{más fuerte} de lo que habíamos afirmado --- coincidencia posición por posición con un ensayo
funcional independiente, en vez de una coincidencia de longitud. Y la contribución mecanística es
\textbf{menor}: ninguna de las coordenadas es nueva.

Lo instructivo es por qué la verificación anterior no lo encontró. Se había corrido una búsqueda web
y no devolvió nada --- pero buscaba el \emph{mecanismo} de la variante, no la cadena de la variante
en texto completo. La conclusión que se sacó entonces --- ``ausencia de resultado no es prueba de
ausencia, así que no afirmar novedad'' --- era la conclusión correcta por la razón equivocada: no
era incertidumbre irreducible, era un paper que nadie había buscado bien. Ese reflejo de retener la
afirmación es lo único que mantuvo el error fuera del registro publicado. \textbf{La consulta
importa tanto como la fuente.}

Una premisa sí sobrevivió a la búsqueda, y ahora está comprobada en vez de asumida: \textbf{no hay
ningún oligonucleótido antisentido publicado contra esta variante}. Corradi et al.\ no mencionan
ninguno. Sangermano et al.~\cite{sangermano2019} --- cuyo título promete corrección de defectos de
splicing de \emph{ABCA4} con AONs, y que era por eso el paper más peligroso para esta premisa ---
diseñaron AONs para otras seis variantes y ninguno para esta. Ninguna de las dos cohortes de 2025
reporta trabajo de AON sobre ella.

% =====================================================================
\section{Los candidatos que sobreviven}
\label{sec:candidatos}

Doce de los 78 candidatos anulan el pseudoexón según la ecuación~\ref{eq:veredicto} en ambos
predictores. Tres son no dominados según la ecuación~\ref{eq:dominancia} (Tabla~\ref{tab:frente}).

\begin{table}[htbp]
\centering
\small
\begin{tabular}{@{}lccccc@{}}
\toprule
\textbf{Candidato} & \textbf{Ventana} & \textbf{Bordes} & $f_1$ & $f_2$ & $f_3$ \\
 & (rel. a la variante) & anulados & bloqueo & off-target & termo \\
\midrule
\texttt{cand\_5992} & $-8$ a $+12$    & donador + aceptor & 1,000 & $-15$ & \phantom{0}8,45 \\
\texttt{cand\_5998} & $-2$ a $+18$    & donador + aceptor & 1,000 & $-16$ & \phantom{0}9,75 \\
\texttt{cand\_5882} & $-118$ a $-98$  & solo aceptor      & 0,994 & $-15$ & 70,75 \\
\bottomrule
\end{tabular}
\caption{El frente de Pareto. Cada candidato es de 20~nt. $f_2$ es el negativo del tramo contiguo
perfecto más largo contra otro transcrito (ecuación~\ref{eq:f2}), así que $-15$ es más seguro que
$-16$. Nótese el compromiso que expone el frente: los dos que mejor bloquean son los \emph{peores}
en propiedades fisicoquímicas, y el de solo-aceptor es al revés. Nada en este proyecto determina qué
compromiso es preferible --- esa es una decisión humana, y por eso el pipeline devuelve tres filas y
no una.}
\label{tab:frente}
\end{table}

\textbf{Ningún candidato fue sintetizado, transfectado ni ensayado.} La Tabla~\ref{tab:frente} es un
orden de prioridad de síntesis, condicionado a que alguien decida sintetizar algo.

% =====================================================================
\section{Lo que este pipeline no hace}

Se declara junto con cada resultado y no es omisible al comunicarlo.

\begin{itemize}[nosep, leftmargin=1.4em]
  \item \textbf{Ninguna validación experimental.} Ningún ASO sintetizado ni ensayado.
  \item \textbf{El proxy de enmascarado no mide eficacia.} La ecuación~\ref{eq:retencion} es binaria
        y total: supone 100\,\% de ocupación y bloqueo perfecto. Necesaria, no suficiente.
  \item \textbf{La termodinámica de PMO es un proxy de otra química}, ahora correctamente orientado
        pero todavía no PMO, y sesgado en una dirección conocida (sección~\ref{sec:termo}).
  \item \textbf{Los umbrales de severidad off-target no están calibrados}: son decisiones de diseño.
  \item \textbf{La concordancia entre predictores no es independiente} y, según el control de tejido,
        en gran medida no informa.
  \item \textbf{El pipeline es ciego al desplazamiento del sitio de splicing}: puntúa cuatro offsets
        fijos y descarta el resto del perfil.
  \item \textbf{El umbral $\tau = 0{,}25$ de la ecuación~\ref{eq:veredicto} es una convención
        declarada}, no un corte medido.
\end{itemize}

% =====================================================================
\section{Discusión}

La contribución científica de este trabajo es modesta y acotada, y un borrador más angosta de lo que
la escribimos primero. El pipeline recupera, desde secuencia sola, un pseudoexón cuyos bordes ya
habían sido establecidos por ensayo funcional en 2023~\cite{corradi2023}; no los descubre. Lo que
queda es un pipeline reproducible que reconstruye un defecto de splicing medido y produce una lista
corta que nadie ha probado.

La contribución metodológica es, creemos, mayor. Un agente de IA produjo un sistema con disciplina de
ingeniería real --- 224 tests, tres entornos documentados, archivos de procedencia, resultados
previos a la corrección archivados --- y simultáneamente produjo conclusiones que no sobrevivieron al
contacto con la revisión adversarial. Y algo central: \textbf{las fallas no fueron aritméticas}. Seis
revisores no encontraron ningún error de cálculo. Fueron fallas de encuadre: qué métrica definir, qué
control no correr, qué afirmación repetir sin comprobar --- y, como muestran las
secciones~\ref{sec:atribucion} y~\ref{sec:novedad}, qué cita no abrir nunca y qué búsqueda no
correr nunca.

Dos de esas fallas aparecieron recién al armar esta lista de referencias, después de que la revisión
adversarial hubiera terminado. Ese momento es en sí mismo un resultado: \textbf{escribir las citas
auditó las premisas del proyecto mejor de lo que lo habían hecho seis revisores especializados},
porque fue el primer paso que obligó a resolver cada afirmación externa contra un identificador y no
contra el recuerdo que el proyecto tenía de ella. Es además el paso más barato de todo el proyecto.

Esa asimetría es el hallazgo que destacaríamos. Es probable que la investigación construida por
agentes sea fuerte donde el trabajo es mecánico y débil donde requiere decidir qué te refutaría. La
revisión adversarial no es entonces control de calidad agregado al final; para este modo de trabajo
es parte del método.

% =====================================================================
\section{Disponibilidad}
\label{sec:disponibilidad}

Código, resultados y comandos de regeneración: \url{https://github.com/smnunez404/bio-oligonucleotidos}
(licencia MIT; datos generados CC BY 4.0). Cada número citado acá tiene un comando de regeneración en
el README del repositorio, y las figuras de este preprint se generan desde \texttt{data/results/} con
un script versionado, así que texto y figura no pueden divergir. El informe de revisión adversarial y
el dossier de afirmaciones se mantienen en el vault de documentación del proyecto y están disponibles
a pedido.

\section*{Agradecimientos}
\addcontentsline{toc}{section}{Agradecimientos}

Agradecemos a los autores de SpliceAI (Illumina), Pangolin (Zeng \& Li), Retina-SpliceAI
(Riepe et al., CMBI/Radboud UMC) y ViennaRNA por publicar su trabajo abiertamente, y a Wang
et al.\ y Kaltak et al.\ por las mediciones experimentales que hicieron posible el anclaje externo.

% =====================================================================
\clearpage
\addcontentsline{toc}{section}{Referencias}
\begin{thebibliography}{99}
\small

\bibitem{wang2025}
Wang Y, Wang P, Yi Z, Ouyang J, Jiang Y, Li S, Jia X, Xiao X, Hejtmancik JF, Sun W, Zhang Q.
\emph{ABCA4 deep intronic variants contributed to nearly half of unsolved Stargardt cases with a
milder phenotype.} Invest Ophthalmol Vis Sci. 2025;66(1):65.
\href{https://doi.org/10.1167/iovs.66.1.65}{doi:10.1167/iovs.66.1.65}

\bibitem{kaltak2023}
Kaltak M, de Bruijn P, Piccolo D, Lee SE, Dulla K, Hoogenboezem T, Beumer W, Webster AR,
Collin RWJ, Cheetham ME, Platenburg G, Swildens J.
\emph{Antisense oligonucleotide therapy corrects splicing in the common Stargardt disease type
1-causing variant ABCA4 c.5461-10T>C.} Mol Ther Nucleic Acids. 2023;31:674--688.
\href{https://doi.org/10.1016/j.omtn.2023.02.020}{doi:10.1016/j.omtn.2023.02.020}

\bibitem{corradi2023}
Corradi Z, Khan M, Hitti-Malin R, Mishra K, Whelan L, Cornelis SS, et al.
\emph{Targeted sequencing and in vitro splice assays shed light on ABCA4-associated retinopathies
missing heritability.} HGG Adv. 2023;4(4):100237.
\href{https://doi.org/10.1016/j.xhgg.2023.100237}{doi:10.1016/j.xhgg.2023.100237}

\bibitem{huang2025}
Huang YS, Lu WT, Chiu IH, Chien PM, Chen YA, Lin CY, Yang CH, Yang CM, Lin CW, Hsu JS, Chen PL,
Chen TC. \emph{Hidden splicing variants in inherited retinal degeneration: discovery and functional
insight.} Invest Ophthalmol Vis Sci. 2025;66(13):12.
\href{https://doi.org/10.1167/iovs.66.13.12}{doi:10.1167/iovs.66.13.12}

\bibitem{sangermano2019}
Sangermano R, Garanto A, Khan M, Runhart EH, Bauwens M, Bax NM, et al.
\emph{Deep-intronic ABCA4 variants explain missing heritability in Stargardt disease and allow
correction of splice defects by antisense oligonucleotides.} Genet Med. 2019;21(8):1751--1760.
\href{https://doi.org/10.1038/s41436-018-0414-9}{doi:10.1038/s41436-018-0414-9}

\bibitem{jaganathan2019}
Jaganathan K, Kyriazopoulou Panagiotopoulou S, McRae JF, et al.
\emph{Predicting splicing from primary sequence with deep learning.}
Cell. 2019;176(3):535--548.e24.
\href{https://doi.org/10.1016/j.cell.2018.12.015}{doi:10.1016/j.cell.2018.12.015}

\bibitem{zeng2022}
Zeng T, Li YI. \emph{Predicting RNA splicing from DNA sequence using Pangolin.}
Genome Biol. 2022;23(1):103.
\href{https://doi.org/10.1186/s13059-022-02664-4}{doi:10.1186/s13059-022-02664-4}

\bibitem{riepe2025}
Riepe TV, de Bruijn SE, Roosing S, Cremers FPM, 't Hoen PAC.
\emph{Deep learning-based prediction of tissue-specific splice sites in the human neural retina.}
bioRxiv, publicado el 13 de febrero de 2025.
\href{https://doi.org/10.1101/2025.02.10.637427}{doi:10.1101/2025.02.10.637427}

\bibitem{sugimoto1995}
Sugimoto N, Nakano S, Katoh M, Matsumura A, Nakamuta H, Ohmichi T, Yoneyama M, Sasaki M.
\emph{Thermodynamic parameters to predict stability of RNA/DNA hybrid duplexes.}
Biochemistry. 1995;34(35):11211--11216.
\href{https://doi.org/10.1021/bi00035a029}{doi:10.1021/bi00035a029}

\bibitem{santalucia1998}
SantaLucia J Jr. \emph{A unified view of polymer, dumbbell, and oligonucleotide DNA
nearest-neighbor thermodynamics.} Proc Natl Acad Sci USA. 1998;95(4):1460--1465.
\href{https://doi.org/10.1073/pnas.95.4.1460}{doi:10.1073/pnas.95.4.1460}

\bibitem{summerton1997}
Summerton J, Weller D. \emph{Morpholino antisense oligomers: design, preparation, and properties.}
Antisense Nucleic Acid Drug Dev. 1997;7(3):187--195.
\href{https://doi.org/10.1089/oli.1.1997.7.187}{doi:10.1089/oli.1.1997.7.187}

\bibitem{maksudov2023}
Maksudov F, Kliuchnikov E, Pierson D, et al.
\emph{Therapeutic phosphorodiamidate morpholino oligonucleotides: physical properties, solution
structures, and folding thermodynamics.} Mol Ther Nucleic Acids. 2023;31:631--647.
\href{https://doi.org/10.1016/j.omtn.2023.02.007}{doi:10.1016/j.omtn.2023.02.007}

\bibitem{park2025}
Park SW, Lee J, Park JW, et al. \emph{Thermodynamic parameter estimation for modified
oligonucleotides using molecular dynamics simulations.} J Phys Chem B. 2025.
\href{https://doi.org/10.1021/acs.jpcb.4c08344}{doi:10.1021/acs.jpcb.4c08344}

\bibitem{lorenz2011}
Lorenz R, Bernhart SH, H\"oner zu Siederdissen C, Tafer H, Flamm C, Stadler PF, Hofacker IL.
\emph{ViennaRNA Package 2.0.} Algorithms Mol Biol. 2011;6:26.
\href{https://doi.org/10.1186/1748-7188-6-26}{doi:10.1186/1748-7188-6-26}

\bibitem{camacho2009}
Camacho C, Coulouris G, Avagyan V, Ma N, Papadopoulos J, Bealer K, Madden TL.
\emph{BLAST+: architecture and applications.} BMC Bioinformatics. 2009;10:421.
\href{https://doi.org/10.1186/1471-2105-10-421}{doi:10.1186/1471-2105-10-421}

\bibitem{cock2009}
Cock PJA, Antao T, Chang JT, et al. \emph{Biopython: freely available Python tools for computational
molecular biology and bioinformatics.} Bioinformatics. 2009;25(11):1422--1423.
\href{https://doi.org/10.1093/bioinformatics/btp163}{doi:10.1093/bioinformatics/btp163}

\bibitem{freeman2018}
Freeman PJ, Hart RK, Gretton LJ, Brookes AJ, Dalgleish R.
\emph{VariantValidator: accurate validation, mapping, and formatting of sequence variation
descriptions.} Hum Mutat. 2018;39(1):61--68.
\href{https://doi.org/10.1002/humu.23348}{doi:10.1002/humu.23348}

\bibitem{scharner2020}
Scharner J, Ma WK, Zhang Q, Lin KT, Rigo F, Bennett CF, Krainer AR.
\emph{Hybridization-mediated off-target effects of splice-switching antisense oligonucleotides.}
Nucleic Acids Res. 2020;48(2):802--816.
\href{https://doi.org/10.1093/nar/gkz1132}{doi:10.1093/nar/gkz1132}

\bibitem{dhuri2020}
Dhuri K, Bechtold C, Quijano E, Pham H, Gupta A, Vikram A, Bahal R.
\emph{Antisense oligonucleotides: an emerging area in drug discovery and development.}
J Clin Med. 2020;9(6):2004.
\href{https://doi.org/10.3390/jcm9062004}{doi:10.3390/jcm9062004}

\bibitem{yoshida2018}
Yoshida T, Naito Y, Sasaki K, Uchida E, Sato Y, Naito M, Kawanishi T, Obika S, Inoue T.
\emph{Estimated number of off-target candidate sites for antisense oligonucleotides in human mRNA
sequences.} Genes Cells. 2018;23(6):448--455.
\href{https://doi.org/10.1111/gtc.12587}{doi:10.1111/gtc.12587}

\bibitem{sullivan2024}
Sullivan PJ, Quinn JMW, Wu W, Pinese M, Cowley MJ.
\emph{SpliceVarDB: a comprehensive database of experimentally validated human splicing variants.}
Am J Hum Genet. 2024;111(10):2164--2175.
\href{https://doi.org/10.1016/j.ajhg.2024.08.002}{doi:10.1016/j.ajhg.2024.08.002}

\end{thebibliography}

\end{document}
